\documentclass[12pt]{article}
\usepackage[utf8]{inputenc}
\usepackage[T1]{fontenc}
\usepackage[margin=1in]{geometry}
\usepackage{setspace}
\usepackage{xcolor}
\usepackage[colorlinks=true, urlcolor=blue!60!black, linkcolor=blue!60!black]{hyperref}
\setlength{\parindent}{2em}
\setlength{\parskip}{1.25em}
\renewcommand{\baselinestretch}{1.0}
\begin{document}
\begin{center}
{\LARGE \textit{Taylor v. United States}, 579 U.S. 301 (2016)\par}
\vspace{0.8em}
{\large Case Brief --- Four-Page Edition\par}
\end{center}
\vspace{1em}

\section{Memory Jogger}

Applying \textit{Gonzales v. Raich} in a 7--1 decision, the Supreme Court held that the government satisfies the Hobbs Act's interstate commerce element merely by proving that a defendant robbed or attempted to rob a drug dealer of drugs or drug proceeds, without proving that the targeted drugs crossed state lines (\href{https://www.law.cornell.edu/supremecourt/text/14-6166}{Cornell LII}).

\section{Detailed Case Facts}

Beginning in 2009, a Roanoke, Virginia gang called the ``Southwest Goonz'' committed home-invasion robberies of drug dealers, who keep cash and drugs at home and rarely call police. David Anthony Taylor took part in two. In August 2009 the gang broke into the home of Josh Whorley, believed to deal ``exotic and high grade'' marijuana, assaulted him and his girlfriend, and left with jewelry, \$40, two cell phones, and one marijuana cigarette. In October 2009 they broke into William Lynch's home, held his wife and children at gunpoint, and left with a cell phone. Taylor was charged under the Hobbs Act, which punishes robbery that ``in any way or degree obstructs, delays, or affects commerce,'' and under \S~924(c) for using a firearm.

\section{Procedural History}

Taylor's first trial in the Western District of Virginia ended in a hung jury after he argued that robbing dealers of local marijuana did not affect interstate commerce. At retrial, the court barred evidence that the dealers sold only local marijuana and ruled that attempting to rob drug dealers satisfied the commerce element as a matter of law. Taylor was convicted and sentenced to 28 years. The Fourth Circuit affirmed, 754 F.3d 217 (2014). The Supreme Court granted certiorari to resolve a split with the Second and Seventh Circuits and, sitting with eight Justices after Justice Scalia's death, affirmed on June 20, 2016.

\section{Judicial Votes}

Affirmed, \textbf{7--1}.
\begin{itemize}
\item \textbf{Majority (7):} Justice Alito, joined by Chief Justice Roberts and Justices Kennedy, Ginsburg, Breyer, Sotomayor, and Kagan.
\item \textbf{Dissent (1):} Justice Thomas. No concurrences.
\end{itemize}

\section{Holding}

In a Hobbs Act prosecution for robbing or attempting to rob a drug dealer, the government need not prove that the drugs crossed or were destined to cross state lines. Proof beyond a reasonable doubt that the defendant knowingly stole or tried to steal drugs or drug proceeds satisfies the commerce element as a matter of law, because under \textit{Raich} the entire market for illegal drugs, including purely intrastate activity, is ``commerce over which the United States has jurisdiction.'' The holding is confined to robberies targeting drug dealers; the Court left open other victims and was not asked to reconsider \textit{Raich}.

\section{Analysis of Opinions}

\textbf{Majority (Alito).} The Hobbs Act's language is ``unmistakably broad'' and reaches the full extent of the commerce power (\textit{Stirone}; \textit{Culbert}). Under \textit{Wickard}, \textit{Lopez}, and \textit{Morrison}, Congress may regulate local economic activity belonging to a class that substantially affects interstate commerce in the aggregate, and \textit{Raich} held that the production and sale of controlled substances, including local marijuana, is such a class. A robber who targets a dealer's drugs or proceeds therefore affects or attempts to affect commerce Congress may regulate. Taylor confused the standard of proof with the meaning of the element: the jury must find the facts beyond a reasonable doubt, but what counts as commerce over which the United States has jurisdiction is a question of law settled by \textit{Raich}, and under \textit{Perez} a minimal effect within a regulated class suffices.

\textbf{Thomas (dissenting).} Under the original meaning, robbery is not ``Commerce,'' and Congress may punish it only when ``necessary and proper'' to regulating interstate commerce, with an ``obvious, simple, and direct relation'' to it---such as robbing an interstate trucker, not a local home invasion. The substantial-effects test has ``no stopping point'' and erases the line between ``what is truly national and what is truly local.'' \textit{Taylor} also expanded \textit{Raich}, which found only a rational basis for congressional judgment. Finally, the decision erodes proof beyond a reasonable doubt of every element (\textit{Winship}), the rule of lenity, and the federalism clear-statement canon (\textit{Jones}).

\section{Examples: Future Applications}

\begin{enumerate}
\item \textbf{Same-Side: Local Methamphetamine Dealer.} A defendant breaks into a rural Ohio dealer's apartment to steal proceeds from methamphetamine made and sold locally. Targeting a dealer's drugs or proceeds satisfies the element as a matter of law.
\item \textbf{Same-Side: Sold-Out Cocaine Seller.} Assailants rob a street-level cocaine seller who has only \$15. As in \textit{Taylor}, an attempt targeting drugs or proceeds suffices even when none are found.
\item \textbf{Opposite-Side: Retired Teacher's Home.} A burglar steals an heirloom watch and \$200 from a home with no drug or commercial connection. \textit{Taylor}'s rule does not apply; the government must prove an actual nexus.
\item \textbf{Opposite-Side: Mugging a Recreational User.} A robber takes a bicycle and wallet from a pedestrian who happens to smoke marijuana. The victim was not targeted as a dealer, so the categorical rule does not apply.
\item \textbf{Fence-Sitter: State-Licensed Dispensary.} A robber holds up a strictly intrastate, state-licensed cannabis dispensary in Colorado, or a homeowner growing wheat for personal use, as Justice Kennedy asked at argument. Whether \textit{Taylor}'s shortcut reaches legal intrastate businesses or other regulated goods is contested.
\end{enumerate}

\section{Critique}

\textbf{Originalist and textualist views.} Echoing Justice Thomas, critics say \textit{Taylor} entrenches \textit{Raich}'s aggregation in criminal law, turning a local home invasion into a 20-year federal felony and letting a rational-basis precedent displace proof of a jurisdictional element. They note the Hobbs Act was aimed at labor-related ``highway robbery'' and extortion obstructing interstate shipments.

\textbf{Progressive and doctrinal views.} Criminal-justice reformers decry federal overcriminalization: Taylor, a ``bit of a tagalong'' who took \$40 and cell phones, received 28 years for conduct Virginia already punished. Scholars question the tension between the claim not to expand \textit{Raich} and using a deference precedent to establish an element of a crime.

\section{Key Quotations}

\begin{quote}
``Rather, to satisfy the Act's commerce element, it is enough that a defendant knowingly stole or attempted to steal drugs or drug proceeds, for, as a matter of law, the market for illegal drugs is `commerce over which the United States has jurisdiction.' ''
\end{quote}
\noindent --- Justice Alito, opinion of the Court, 579 U.S. at 309.

\begin{quote}
``Our holding today is limited to cases in which the defendant targets drug dealers for the purpose of stealing drugs or drug proceeds.''
\end{quote}
\noindent --- Justice Alito, opinion of the Court, 579 U.S. at 310.

\begin{quote}
``By applying the substantial-effects test to the criminal prohibition before us, the Court effectively gives Congress a police power.''
\end{quote}
\noindent --- Justice Thomas, dissenting, 579 U.S. at 319.

\end{document}
